\section{Justificativa}
\label{justificativa}

A Justificativa apresenta as razões teóricas, práticas e institucionais que tornam a realização do trabalho oportuna e relevante. Nesta seção, o autor deve demonstrar o impacto esperado da solução proposta perante a comunidade acadêmica e a sociedade.

Em trabalhos desenvolvidos na CESAR School, a justificativa deve ser sustentada sob três dimensões complementares:

\begin{itemize}
    \item \textbf{Relevância Científica e Acadêmica:} Como o trabalho contribui para o estado da arte, preenchendo lacunas conceituais, propondo novas taxonomias ou fornecendo dados empíricos inéditos.
    \item \textbf{Relevância Tecnológica e de Mercado:} Como a solução resolve problemas reais da indústria de software, design ou gestão, promovendo ganhos de eficiência, redução de custos ou automação de processos.
    \item \textbf{Relevância Social e Institucional:} Como os resultados geram valor para a sociedade, alinhando-se aos valores de inovação aberta, sustentabilidade e transformação digital.
\end{itemize}

A fundamentação da justificativa deve ser acompanhada de evidências concretas, tais como relatórios de mercado, estatísticas de uso de sistemas ou referências da literatura especializada.
